\documentclass[a4paper,11pt]{article}

% ---- Encoding & fonts -------------------------------------------------------
\usepackage[T1]{fontenc}
\usepackage[utf8]{inputenc}
\usepackage{lmodern}
\usepackage{microtype}              % protrusion/kerning -> fewer overfull boxes

% ---- Geometry (~1in margins) ------------------------------------------------
\usepackage[a4paper,margin=1in,headheight=15pt,headsep=14pt]{geometry}

% ---- Math -------------------------------------------------------------------
\usepackage{amsmath,amssymb}

% ---- Graphics, tables, colour ----------------------------------------------
\usepackage{graphicx}
\usepackage{booktabs}
\usepackage[table,svgnames]{xcolor}
\usepackage{adjustbox}

% ---- Callout boxes (breakable pill-tab tcolorboxes) -------------------------
\usepackage[most]{tcolorbox}
\tcbuselibrary{skins,breakable}

\usepackage{pifont}
\IfFileExists{bbding.sty}{%
  \usepackage{bbding}%
  \newcommand{\glyphHand}{\Pointinghand}%
  \newcommand{\glyphPencil}{\Pencil}%
}{%
  \newcommand{\glyphHand}{\ding{43}}%
  \newcommand{\glyphPencil}{\ding{46}}%
}

% ---- Lists & spacing --------------------------------------------------------
\usepackage{enumitem}
\setlist{leftmargin=1.5em,topsep=4pt,itemsep=3pt}

\newlist{steps}{enumerate}{1}
\setlist[steps]{label=\textbf{Step \arabic*.},
  leftmargin=4.4em, labelwidth=3.8em, labelsep=0.6em, labelindent=0pt,
  align=left, topsep=5pt, itemsep=6pt}
\usepackage{parskip}

% ---- Headers / footers ------------------------------------------------------
\usepackage{fancyhdr}

% ---- Section styling --------------------------------------------------------
\usepackage{titlesec}

% ---- Links (hidden) + URL line-breaking so long URLs never overflow ---------
\usepackage[hidelinks]{hyperref}
\usepackage{xurl}
\hypersetup{breaklinks=true}

% =============================================================================
%  COLOURS — navy chrome + the FIVE taxonomy ramps
% =============================================================================
\definecolor{navy}{HTML}{21355E}        % banner + headings
\definecolor{navyrule}{HTML}{21355E}
\definecolor{inktext}{HTML}{1A202C}     % body ink
\definecolor{mutedtext}{HTML}{6B7280}   % header / meta grey

% Intuition — blue
\definecolor{intuFrame}{HTML}{2C5AA0}   \definecolor{intuFill}{HTML}{F4F7FC}
% Everyday — amber
\definecolor{everyFrame}{HTML}{8A5A1E}  \definecolor{everyFill}{HTML}{FBF1E3}
% Worked — green
\definecolor{workFrame}{HTML}{2E7D52}   \definecolor{workFill}{HTML}{F1F8F3}
% Watch out — red
\definecolor{watchFrame}{HTML}{B23A48}  \definecolor{watchFill}{HTML}{FBF1F2}
% Key takeaway — purple/violet
\definecolor{keyFrame}{HTML}{6A4C93}    \definecolor{keyFill}{HTML}{F4F0F9}
% Wait, what? — teal
\definecolor{waitFrame}{HTML}{0F766E}   \definecolor{waitFill}{HTML}{EEF7F5}

\color{inktext}
\hypersetup{linkcolor=navy,urlcolor=navy,citecolor=navy}

\newcommand{\CourseShort}{DNN Companion}
\newcommand{\SessionNo}{11}
\newcommand{\LectureTitle}{Seq2Seq Architectures \& Attention Mechanisms}

\pagestyle{fancy}
\fancyhf{}
\fancyhead[L]{\footnotesize\color{mutedtext}\textsf{\CourseShort}}
\fancyhead[R]{\footnotesize\color{mutedtext}\textsf{Session \SessionNo\ \textperiodcentered\ \LectureTitle}}
\fancyfoot[L]{\footnotesize\color{mutedtext}\textsf{WILP, BITS Pilani}}
\fancyfoot[C]{\footnotesize\color{mutedtext}\thepage}
\renewcommand{\headrulewidth}{0.4pt}
\renewcommand{\footrulewidth}{0pt}
\renewcommand{\headrule}{\hbox to\headwidth{\color{navyrule!55}\leaders\hrule height \headrulewidth\hfill}}

\titleformat{\section}
  {\normalfont\Large\bfseries\sffamily\color{navy}}
  {\thesection}{0.8em}{}
  [\vspace{2pt}\color{navyrule!40}\hrule height 0.6pt]
\titlespacing*{\section}{0pt}{18pt}{8pt}

\titleformat{\subsection}
  {\normalfont\large\bfseries\sffamily\color{navy}}
  {\thesubsection}{0.6em}{}
\titlespacing*{\subsection}{0pt}{12pt}{4pt}

\newcommand{\secsub}[1]{%
  \vspace{-6pt}{\small\sffamily\color{mutedtext}\textbf{#1}}\par\vspace{8pt}%
}

% =============================================================================
%  TCOLORBOX CALLOUT MACROS
% =============================================================================
\newtcolorbox{intuition}[1][Mathematical Intuition]{%
  enhanced, breakable, frame hidden, interior hidden,
  toptitle=5pt, bottomtitle=5pt, top=14pt, bottom=10pt, left=12pt, right=12pt,
  colback=intuFill, colframe=intuFrame,
  overlay unbroken and first={
    \draw[intuFrame, lw=1pt] (frame.south west) rectangle (frame.north east);
    \node[anchor=west, fill=intuFrame, text=white, font=\sffamily\bfseries\small,
          inner xsep=10pt, inner ysep=4pt, rounded corners=3pt]
      at ([xshift=12pt, yshift=0pt]frame.north west) {\ding{168}\ \ #1};
  }%
}

\newtcolorbox{everyday}[1][Everyday Picture]{%
  enhanced, breakable, frame hidden, interior hidden,
  toptitle=5pt, bottomtitle=5pt, top=14pt, bottom=10pt, left=12pt, right=12pt,
  colback=everyFill, colframe=everyFrame,
  overlay unbroken and first={
    \draw[everyFrame, lw=1pt] (frame.south west) rectangle (frame.north east);
    \node[anchor=west, fill=everyFrame, text=white, font=\sffamily\bfseries\small,
          inner xsep=10pt, inner ysep=4pt, rounded corners=3pt]
      at ([xshift=12pt, yshift=0pt]frame.north west) {\glyphHand\ \ #1};
  }%
}

\newtcolorbox{worked}[1][Worked Example]{%
  enhanced, breakable, frame hidden, interior hidden,
  toptitle=5pt, bottomtitle=5pt, top=14pt, bottom=10pt, left=12pt, right=12pt,
  colback=workFill, colframe=workFrame,
  overlay unbroken and first={
    \draw[workFrame, lw=1pt] (frame.south west) rectangle (frame.north east);
    \node[anchor=west, fill=workFrame, text=white, font=\sffamily\bfseries\small,
          inner xsep=10pt, inner ysep=4pt, rounded corners=3pt]
      at ([xshift=12pt, yshift=0pt]frame.north west) {\glyphPencil\ \ #1};
  }%
}

\newtcolorbox{watchout}[1][Watch Out!]{%
  enhanced, breakable, frame hidden, interior hidden,
  toptitle=5pt, bottomtitle=5pt, top=14pt, bottom=10pt, left=12pt, right=12pt,
  colback=watchFill, colframe=watchFrame,
  overlay unbroken and first={
    \draw[watchFrame, lw=1pt] (frame.south west) rectangle (frame.north east);
    \node[anchor=west, fill=watchFrame, text=white, font=\sffamily\bfseries\small,
          inner xsep=10pt, inner ysep=4pt, rounded corners=3pt]
      at ([xshift=12pt, yshift=0pt]frame.north west) {\ding{55}\ \ #1};
  }%
}

\newtcolorbox{keytake}[1][Key Takeaway]{%
  enhanced, breakable, frame hidden, interior hidden,
  toptitle=5pt, bottomtitle=5pt, top=14pt, bottom=10pt, left=12pt, right=12pt,
  colback=keyFill, colframe=keyFrame,
  overlay unbroken and first={
    \draw[keyFrame, lw=1pt] (frame.south west) rectangle (frame.north east);
    \node[anchor=west, fill=keyFrame, text=white, font=\sffamily\bfseries\small,
          inner xsep=10pt, inner ysep=4pt, rounded corners=3pt]
      at ([xshift=12pt, yshift=0pt]frame.north west) {\ding{51}\ \ #1};
  }%
}

\newtcolorbox{waitwhat}[1][Wait, What?]{%
  enhanced, breakable, frame hidden, interior hidden,
  toptitle=5pt, bottomtitle=5pt, top=14pt, bottom=10pt, left=12pt, right=12pt,
  colback=waitFill, colframe=waitFrame,
  overlay unbroken and first={
    \draw[waitFrame, lw=1pt] (frame.south west) rectangle (frame.north east);
    \node[anchor=west, fill=waitFrame, text=white, font=\sffamily\bfseries\small,
          inner xsep=10pt, inner ysep=4pt, rounded corners=3pt]
      at ([xshift=12pt, yshift=0pt]frame.north west) {\small!?\ \ #1};
  }%
}

\newcommand{\housefig}[2]{%
  \begin{figure}[htbp]
    \centering
    \includegraphics[width=0.92\linewidth]{#1}
    \caption{#2}
  \end{figure}
}

\newcommand{\flipanswer}[1]{%
  \begin{adjustbox}{angle=180}%
    \begin{minipage}{\linewidth}%
      \small\color{mutedtext}#1%
    \end{minipage}%
  \end{adjustbox}%
}

\begin{document}

% Banner Header
\begin{tcolorbox}[
  enhanced,
  colback=navy,
  colframe=navy,
  arc=0pt, outer arc=0pt,
  left=16pt, right=16pt, top=14pt, bottom=14pt
]
  \color{white}
  \begin{minipage}{0.72\linewidth}
    {\fontfamily{phv}\selectfont
      {\small\bfseries\scshape WILP, BITS Pilani \hfill Deep Neural Networks (AIML ZG511)}\par
      \vspace{4pt}
      {\LARGE\bfseries Session 11: Seq2Seq Architectures \& Attention Mechanisms \par}
      \vspace{4pt}
      {\small Encoder-Decoder Bottlenecks, Additive vs Multiplicative Attention, and Soft Alignment\par}
    }
  \end{minipage}%
  \hfill
  \begin{minipage}{0.24\linewidth}
    \centering
    \color{white!80}\sffamily\scriptsize
    \textbf{Module 8}\\[2pt]
    Session 11 Reader\\[2pt]
    Full Worked Solutions
  \end{minipage}
\end{tcolorbox}

\vspace{10pt}

% Section 1
\section{Sequence-to-Sequence (Seq2Seq) \& The Bottleneck Problem}
\secsub{Seq2Seq networks encode a whole input sequence into a single vector, causing a major memory bottleneck.}

Standard sequence-to-sequence (Seq2Seq) architectures consist of an Encoder RNN and a Decoder RNN. The Encoder reads input sequence $(x_1, x_2, \dots, x_T)$ and compresses it into a single context vector $c = h_T$. The Decoder generates target sequence $(y_1, y_2, \dots, y_{T'})$ from $c$.

\housefig{figures/seq2seq_bottleneck.png}{The Seq2Seq Encoder-Decoder bottleneck problem. All input context is squeezed into a single fixed vector $c$.}

\begin{everyday}[The Interpreter's Limit]
Imagine listening to a 20-minute speech in French, taking no notes, and then translating the whole speech into English strictly from memory. Long sentences get forgotten!
\end{everyday}

% Section 2
\section{The Attention Mechanism: Soft Alignment}
\secsub{Attention lets the decoder look back at all encoder hidden states dynamically at every step.}

Instead of forcing a single context vector $c$, Attention computes a unique context vector $c_i$ for each decoder step $i$ as a weighted sum:
$$c_i = \sum_{j=1}^{T} \alpha_{ij} h_j$$
where $\alpha_{ij}$ represents the alignment weight between decoder state $s_{i-1}$ and encoder state $h_j$.

\housefig{figures/attention_heatmap.png}{Attention alignment matrix $\alpha_{ij}$ showing soft word alignments between source and target sentences.}

% Section 3
\section{Computing Attention Weights: Step-by-Step Math}
\secsub{Alignment scores are converted to a probability distribution over encoder states using softmax.}

\begin{intuition}[Alignment Score Formulations]
1. \textbf{Bahdanau (Additive) Attention:}
$$e_{ij} = v_a^T \tanh(W_a s_{i-1} + U_a h_j)$$
2. \textbf{Luong (Multiplicative) Attention:}
$$e_{ij} = s_{i-1}^T W_a h_j \quad \text{or simple dot product } e_{ij} = s_{i-1}^T h_j$$
\end{intuition}

\begin{worked}[Full Worked Attention Numerical]
Suppose an encoder has $T=3$ hidden states:
$$h_1 = \begin{bmatrix} 1 \\ 0 \end{bmatrix}, \quad h_2 = \begin{bmatrix} 0 \\ 2 \end{bmatrix}, \quad h_3 = \begin{bmatrix} 1 \\ 1 \end{bmatrix}$$
The decoder state is $s_{i-1} = \begin{bmatrix} 2 \\ 1 \end{bmatrix}$. We use Dot-Product Attention $e_{ij} = s_{i-1}^T h_j$.

\begin{steps}
\item Compute raw alignment scores $e_{ij}$:
$$e_{i1} = \begin{bmatrix} 2 & 1 \end{bmatrix} \begin{bmatrix} 1 \\ 0 \end{bmatrix} = 2.0, \quad e_{i2} = \begin{bmatrix} 2 & 1 \end{bmatrix} \begin{bmatrix} 0 \\ 2 \end{bmatrix} = 2.0$$
$$e_{i3} = \begin{bmatrix} 2 & 1 \end{bmatrix} \begin{bmatrix} 1 \\ 1 \end{bmatrix} = 3.0$$

\item Apply Softmax to compute attention weights $\alpha_{ij}$:
$$\exp(e_{i1}) = 7.389, \quad \exp(e_{i2}) = 7.389, \quad \exp(e_{i3}) = 20.086$$
$$\sum_{j=1}^{3} \exp(e_{ij}) = 7.389 + 7.389 + 20.086 = 34.864$$
$$\alpha_{i1} = 0.2119, \quad \alpha_{i2} = 0.2119, \quad \alpha_{i3} = 0.5761$$

\item Calculate dynamic context vector $c_i$:
\begin{align*}
c_i &= 0.2119 \begin{bmatrix} 1 \\ 0 \end{bmatrix} + 0.2119 \begin{bmatrix} 0 \\ 2 \end{bmatrix} + 0.5761 \begin{bmatrix} 1 \\ 1 \end{bmatrix} \\
&= \begin{bmatrix} 0.2119 \\ 0 \end{bmatrix} + \begin{bmatrix} 0 \\ 0.4238 \end{bmatrix} + \begin{bmatrix} 0.5761 \\ 0.5761 \end{bmatrix} = \begin{bmatrix} 0.7880 \\ 0.9999 \end{bmatrix}
\end{align*}
\end{steps}
\end{worked}

\begin{keytake}[Additive vs Multiplicative Attention]
\begin{center}
\begin{tabular}{lll}
\toprule
\textbf{Property} & \textbf{Bahdanau (Additive)} & \textbf{Luong (Multiplicative)} \\
\midrule
Score Formula & $v_a^T \tanh(W_a s + U_a h)$ & $s^T W_a h$ \\
Parameters & Requires $W_a, U_a, v_a$ & Requires only $W_a$ (or 0 for dot) \\
Complexity & Higher & Lower (optimized BLAS dot products) \\
\bottomrule
\end{tabular}
\end{center}
\end{keytake}

% Section 4
\section{Self-Test \& Pocket Recap}
\secsub{Test your knowledge on attention mechanisms and context computation.}

\subsection*{Self-Test Questions}
1. \textbf{Question:} Why does Luong attention scale better computationally than Bahdanau attention?\\
\flipanswer{\textbf{Answer:} Luong attention uses matrix dot products ($s^T W_a h$) which map directly to highly optimized GPU matrix multiplication routines, whereas Bahdanau attention requires non-linear activation functions ($\tanh$).}

\vspace{10pt}

2. \textbf{Question:} If all alignment scores $e_{ij}$ are equal to $5.0$ for sequence length $T=4$, what are the attention weights?\\
\flipanswer{\textbf{Answer:} Since all raw scores are equal, softmax produces a uniform distribution: $\alpha_{ij} = \frac{1}{4} = 0.25$ for all $j=1,2,3,4$.}

\end{document}
